\documentclass[11pt]{article}

\usepackage[margin=1in]{geometry}
\usepackage{booktabs}
\usepackage{graphicx}
\usepackage{parskip}
\usepackage{enumitem}
\usepackage{hyperref}
\usepackage{array}

\title{\textbf{CLAHE Preprocessing Experiment Report}\\[4pt]
Query-Side Contrast Enhancement for Nocturnal VPR}
\author{Rhoda Ojetola \and Peter Adeyemo \and Samuel Olusola}
\date{April 25, 2026}

\begin{document}
\maketitle

\section*{Purpose}
This experiment studies whether \textbf{CLAHE preprocessing} helps nocturnal visual place recognition in our project. The idea is simple: instead of changing the descriptors themselves, we preprocess the \textbf{query/night images} before feature extraction and compare the new runs against the existing no-preprocessing baselines.

\section*{Hypothesis}
If campus night images are genuinely dark and low-contrast, then improving local contrast on the query side should reduce the day-night appearance gap and improve retrieval. We therefore expect CLAHE to help \textbf{more on campus} than on GardensPoint.

\section*{What Was Implemented}
We added a shared preprocessing module in:
\begin{itemize}[leftmargin=*]
  \item \texttt{evaluation/preprocessing.py}
\end{itemize}

The shared evaluation scripts now support:
\begin{itemize}[leftmargin=*]
  \item \texttt{--preprocess none}
  \item \texttt{--preprocess clahe\_query}
  \item \texttt{--preprocess clahe\_all}
  \item \texttt{--clahe\_clip\_limit}
  \item \texttt{--clahe\_tile\_grid}
\end{itemize}

For this report we used \texttt{clahe\_query}, because that matches the intended use case: improve the night/query images without altering the day map.

\section*{Experimental Setup}
\textbf{CLAHE settings used in all runs}
\begin{itemize}[leftmargin=*]
  \item mode: \texttt{clahe\_query}
  \item clip limit: \texttt{2.0}
  \item tile grid size: \texttt{8}
\end{itemize}

\textbf{Descriptors tested}
\begin{itemize}[leftmargin=*]
  \item \textbf{CosPlace}: strong VPR-specific baseline already used throughout the project
  \item \textbf{EigenPlaces}: another strong modern baseline that performs well on ranking
  \item \textbf{SALAD}: strongest practical modern descriptor in our recent experiments
  \item \textbf{SelaVPR++}: strongest benchmark descriptor overall in our recent experiments
\end{itemize}

\textbf{Datasets / protocols tested}
\begin{itemize}[leftmargin=*]
  \item GardensPoint benchmark
  \item campus strict image-level protocol
  \item campus place-level protocol
\end{itemize}

\section*{Commands Used}
\begin{verbatim}
python demo.py --descriptor SALAD --dataset GardensPoint \
  --preprocess clahe_query --save_results

python test_campus_dataset.py --descriptor SALAD \
  --preprocess clahe_query --save_results

python test_campus_dataset_place_level.py --descriptor SALAD \
  --preprocess clahe_query --save_results

python demo.py --descriptor "SelaVPR++" --dataset GardensPoint \
  --preprocess clahe_query --save_results

python test_campus_dataset.py --descriptor "SelaVPR++" \
  --preprocess clahe_query --save_results

python test_campus_dataset_place_level.py --descriptor "SelaVPR++" \
  --preprocess clahe_query --save_results

python demo.py --descriptor CosPlace --dataset GardensPoint \
  --preprocess clahe_query --save_results

python test_campus_dataset.py --descriptor CosPlace \
  --preprocess clahe_query --save_results

python test_campus_dataset_place_level.py --descriptor CosPlace \
  --preprocess clahe_query --save_results

python demo.py --descriptor EigenPlaces --dataset GardensPoint \
  --preprocess clahe_query --save_results

python test_campus_dataset.py --descriptor EigenPlaces \
  --preprocess clahe_query --save_results

python test_campus_dataset_place_level.py --descriptor EigenPlaces \
  --preprocess clahe_query --save_results
\end{verbatim}

\section*{Results}
\subsection*{CosPlace}
\begin{center}
\begin{tabular}{>{\raggedright\arraybackslash}p{3.0cm}cccccc}
\toprule
\textbf{Setting} & \textbf{AUC} & \textbf{$\Delta$AUC} & \textbf{R@100P} & \textbf{$\Delta$R@100P} & \textbf{R@1} & \textbf{$\Delta$R@1} \\
\midrule
GardensPoint baseline & 0.592 & -- & 0.270 & -- & 0.350 & -- \\
GardensPoint + CLAHE & 0.573 & -0.019 & 0.225 & -0.045 & 0.330 & -0.020 \\
\midrule
Campus strict baseline & 0.470 & -- & 0.030 & -- & 0.727 & -- \\
Campus strict + CLAHE & 0.454 & -0.016 & 0.061 & +0.031 & 0.727 & +0.000 \\
\midrule
Campus place baseline & 0.627 & -- & 0.114 & -- & 0.969 & -- \\
Campus place + CLAHE & 0.629 & +0.002 & 0.110 & -0.004 & 0.969 & +0.000 \\
\bottomrule
\end{tabular}
\end{center}

\subsection*{EigenPlaces}
\begin{center}
\begin{tabular}{>{\raggedright\arraybackslash}p{3.0cm}cccccc}
\toprule
\textbf{Setting} & \textbf{AUC} & \textbf{$\Delta$AUC} & \textbf{R@100P} & \textbf{$\Delta$R@100P} & \textbf{R@1} & \textbf{$\Delta$R@1} \\
\midrule
GardensPoint baseline & 0.675 & -- & 0.220 & -- & 0.365 & -- \\
GardensPoint + CLAHE & 0.655 & -0.020 & 0.240 & +0.020 & 0.355 & -0.010 \\
\midrule
Campus strict baseline & 0.443 & -- & 0.000 & -- & 0.788 & -- \\
Campus strict + CLAHE & 0.488 & +0.045 & 0.030 & +0.030 & 0.879 & +0.091 \\
\midrule
Campus place baseline & 0.663 & -- & 0.110 & -- & 0.922 & -- \\
Campus place + CLAHE & 0.663 & +0.000 & 0.100 & -0.010 & 0.938 & +0.016 \\
\bottomrule
\end{tabular}
\end{center}

\subsection*{SALAD}
\begin{center}
\begin{tabular}{>{\raggedright\arraybackslash}p{3.0cm}cccccc}
\toprule
\textbf{Setting} & \textbf{AUC} & \textbf{$\Delta$AUC} & \textbf{R@100P} & \textbf{$\Delta$R@100P} & \textbf{R@1} & \textbf{$\Delta$R@1} \\
\midrule
GardensPoint baseline & 0.924 & -- & 0.385 & -- & 0.450 & -- \\
GardensPoint + CLAHE & 0.896 & -0.028 & 0.430 & +0.045 & 0.415 & -0.035 \\
\midrule
Campus strict baseline & 0.504 & -- & 0.030 & -- & 0.818 & -- \\
Campus strict + CLAHE & 0.522 & +0.018 & 0.030 & +0.000 & 0.848 & +0.030 \\
\midrule
Campus place baseline & 0.653 & -- & 0.130 & -- & 0.969 & -- \\
Campus place + CLAHE & 0.651 & -0.002 & 0.128 & -0.002 & 0.969 & +0.000 \\
\bottomrule
\end{tabular}
\end{center}

\subsection*{SelaVPR++}
\begin{center}
\begin{tabular}{>{\raggedright\arraybackslash}p{3.0cm}cccccc}
\toprule
\textbf{Setting} & \textbf{AUC} & \textbf{$\Delta$AUC} & \textbf{R@100P} & \textbf{$\Delta$R@100P} & \textbf{R@1} & \textbf{$\Delta$R@1} \\
\midrule
GardensPoint baseline & 0.937 & -- & 0.355 & -- & 0.465 & -- \\
GardensPoint + CLAHE & 0.933 & -0.004 & 0.495 & +0.140 & 0.440 & -0.025 \\
\midrule
Campus strict baseline & 0.464 & -- & 0.061 & -- & 0.909 & -- \\
Campus strict + CLAHE & 0.481 & +0.017 & 0.061 & +0.000 & 0.909 & +0.000 \\
\midrule
Campus place baseline & 0.658 & -- & 0.158 & -- & 0.969 & -- \\
Campus place + CLAHE & 0.659 & +0.001 & 0.152 & -0.006 & 0.969 & +0.000 \\
\bottomrule
\end{tabular}
\end{center}

\section*{Interpretation}
\subsection*{1. CLAHE does \emph{not} help uniformly}
The effect is dataset-dependent and metric-dependent.
\begin{itemize}[leftmargin=*]
  \item On \textbf{GardensPoint}, CLAHE reduced AUC for \emph{all four} descriptors and reduced R@1 for all four as well.
  \item On \textbf{campus strict}, CLAHE helped the stronger learned descriptors most clearly: EigenPlaces, SALAD, and SelaVPR++ all improved AUC, while EigenPlaces and SALAD also improved R@1.
  \item On \textbf{campus place-level}, CLAHE made only small changes overall, because the task is already less brittle once multiple valid views are credited.
\end{itemize}

\subsection*{2. This matches the EDA}
Our raw-data EDA already suggested this outcome:
\begin{itemize}[leftmargin=*]
  \item GardensPoint night images appear \textbf{exposure-compensated / over-brightened}.
  \item Campus night images are \textbf{genuinely darker} and slightly lower contrast.
\end{itemize}

So CLAHE is more likely to help on campus, where there is a real contrast problem, than on GardensPoint, where the night images are already brightened.

\subsection*{3. Why did R@100P improve on GardensPoint?}
For some descriptors, CLAHE improved \textbf{R@100P} on GardensPoint even though AUC and R@1 went down.

This suggests the following interpretation:
\begin{itemize}[leftmargin=*]
  \item CLAHE may slightly clean up the \emph{thresholded} match distribution, making it easier to reach a perfect-precision operating point.
  \item But it also distorts some local appearance cues enough to hurt top-1 ranking.
\end{itemize}

So CLAHE may help the binary decision boundary in some cases while hurting the very best-match ordering.

\subsection*{4. Why is campus place-level almost unchanged?}
The place-level protocol already gives credit for multiple valid views of the same place. That makes the task less brittle than the strict image-level protocol. Because of that, CLAHE has less room to help: the descriptor is already solving most of the task once multi-view positives are allowed.

\subsection*{5. Best practical conclusion}
The strongest practical conclusion from these runs is:
\begin{itemize}[leftmargin=*]
  \item \textbf{Do not treat CLAHE as a universal default.}
  \item It is \textbf{most promising for the strict campus setup}, where the query images are truly dark and the exact image-level matching protocol is sensitive to appearance shifts.
  \item It is \textbf{not clearly useful on GardensPoint}, and may actually hurt top-1 retrieval there.
\end{itemize}

\subsection*{6. Descriptor-specific summary}
\begin{itemize}[leftmargin=*]
  \item \textbf{CosPlace}: CLAHE does not meaningfully help. It slightly hurts GardensPoint and strict-campus AUC, while leaving R@1 unchanged on campus.
  \item \textbf{EigenPlaces}: this is the clearest positive baseline result. CLAHE noticeably improves strict-campus AUC and R@1, while leaving place-level AUC unchanged.
  \item \textbf{SALAD}: CLAHE helps the strict campus setup, but hurts GardensPoint overall.
  \item \textbf{SelaVPR++}: CLAHE gives a small strict-campus AUC gain, but mostly leaves the already-strong place-level ranking unchanged.
\end{itemize}

\section*{Qualitative Note}
The strict campus CLAHE runs reduced the number of false positives while keeping or improving top-k retrieval:
\begin{itemize}[leftmargin=*]
  \item EigenPlaces strict: AUC rose from 0.443 to 0.488 and R@1 rose from 0.788 to 0.879
  \item SALAD strict: FP fell from 6 to 4, and R@1 rose from 0.818 to 0.848
  \item SelaVPR++ strict: FP stayed at 2, AUC rose from 0.464 to 0.481, while R@1 stayed at 0.909
\end{itemize}

This is a good sign that CLAHE may be more useful as a \textbf{campus-specific query enhancement} than as a general benchmark improvement.

\section*{Run Artifacts}
The new CLAHE runs are saved under:
\begin{itemize}[leftmargin=*]
  \item \texttt{output\_images/runs/benchmarks/0015\_demo\_gardenspoint\_salad\_clahe\_query\_...}
  \item \texttt{output\_images/runs/campus\_strict/0008\_campus\_strict\_salad\_clahe\_query\_...}
  \item \texttt{output\_images/runs/campus\_place\_level/0007\_campus\_place\_level\_salad\_clahe\_query\_...}
  \item \texttt{output\_images/runs/benchmarks/0016\_demo\_gardenspoint\_selavpr\_clahe\_query\_...}
  \item \texttt{output\_images/runs/campus\_strict/0009\_campus\_strict\_selavpr\_clahe\_query\_...}
  \item \texttt{output\_images/runs/campus\_place\_level/0008\_campus\_place\_level\_selavpr\_clahe\_query\_...}
  \item \texttt{output\_images/runs/benchmarks/0017\_demo\_gardenspoint\_cosplace\_clahe\_query\_...}
  \item \texttt{output\_images/runs/campus\_strict/0010\_campus\_strict\_cosplace\_clahe\_query\_...}
  \item \texttt{output\_images/runs/campus\_place\_level/0009\_campus\_place\_level\_cosplace\_clahe\_query\_...}
  \item \texttt{output\_images/runs/benchmarks/0018\_demo\_gardenspoint\_eigenplaces\_clahe\_query\_...}
  \item \texttt{output\_images/runs/campus\_strict/0011\_campus\_strict\_eigenplaces\_clahe\_query\_...}
  \item \texttt{output\_images/runs/campus\_place\_level/0010\_campus\_place\_level\_eigenplaces\_clahe\_query\_...}
\end{itemize}

\section*{Conclusion}
CLAHE is worth keeping in the project as a configurable option, but the experiments show that its usefulness is \textbf{domain-specific}. It helps most where the night images are truly dark and contrast-limited, which is exactly what our campus EDA suggested. The cleanest next step is therefore to keep CLAHE as an optional preprocessing branch for \textbf{campus-style deployment}, not as a universal preprocessing default for every benchmark.

\end{document}
